\section{Ablation}
\label{sec:ablation}

The protocol has components of two kinds: the rules on the row sets that the split enforces, and the gates that test the rules. Two ablations separate what each contributes. The first removes one rule at a time and reads the cost in the bias and the coverage of the resulting leak. The second keeps every subset of the gates and reads what the subset detects.

\subsection{Removing one rule of the split}

Table~\ref{tab:ablation-components} gives, for each rule that is removed, the worst of the leaks that the removal allows. The rule that the statistics are fitted on the training slice only is the most valuable: its removal in the whole fitting path adds \nv{syn.w.fit.bias}$\,\pm\,$\nv{syn.w.fit.ci} to the AUC and lowers the coverage of the interval to \nv{syn.w.fit.cov}\%. The value is not spread evenly over the steps. Of the five steps whose removal is covered by this rule, two cost nothing measurable on the planted data (imputation, \nv{syn.w.imp.bias}, and bin edges, \nv{syn.w.bin.bias}); the screen costs \nv{syn.w.scr.bias}; and the weight of evidence costs \nv{syn.w.woe.bias}, which is more than two thirds of the cost of reading the holdout in every step, so that the target-aware transform is where a scorecard is most exposed. The rule that the configuration is chosen inside the training slice is the cheapest: its removal costs \nv{syn.w.sel.bias}$\,\pm\,$\nv{syn.w.sel.ci} with eight configurations and nothing measurable with two, which is the behaviour of Proposition~\ref{prop:selection} for correlated configurations. The rule that the calibration slice is disjoint from the test slice costs \nv{syn.w.caltest.bias}$\,\pm\,$\nv{syn.w.caltest.ci} in AUC, lowers the coverage to \nv{syn.w.caltest.cov}\% and is the only rule whose removal makes the calibration error look better than it is. The rule that the cross-validation folds contain the preprocessing is the most expensive when a scorecard is validated by cross-validation rather than by a holdout, with \nv{syn.w.cvpre.bias}$\,\pm\,$\nv{syn.w.cvpre.ci}; on the small sample it reaches \nv{small.w.cvpre.bias}. No rule is redundant for the estimate: each removal produces a bias whose interval excludes zero on the planted data, except the two input-only steps, and those are the ones that cost nothing in bias while a gate still fails on them.

\begin{table}[htbp]
\centering
\small
\caption{Removing one component of the protocol at a time on the planted data of the main comparison (WoE scorecard): the leak that results from the removal with the largest AUC bias, its bias with the 95\% half-width, the coverage of the 95\% interval and the Brier bias. The reference row of the protocol is in the main comparison table.}
\label{tab:ablation-components}
\fitwidth{%
\begin{tabular}{llrrr}
\toprule
Component of the protocol removed & Worst resulting leak & AUC bias & Coverage \% & Brier bias \\
\midrule
Statistics fitted on the training slice only & everything fitted on training and holdout & $+0.025\pm0.002$ & 66.7 & $-$0.0088 \\
Configuration chosen inside the training slice & configuration chosen on the holdout & $+0.005\pm0.002$ & 96.0 & $-$0.0018 \\
Calibration slice disjoint from the test slice & calibration map fitted on the holdout & $+0.016\pm0.002$ & 88.0 & $-$0.0081 \\
Preprocessing inside the cross-validation folds & five-fold, preprocessing before the folds & $+0.054\pm0.001$ & -- & -- \\
\bottomrule
\end{tabular}%
}
\end{table}


\subsection{Keeping a subset of the gates}

Table~\ref{tab:ablation-gates} gives the detection of every subset of the three gates over the \nv{h5.total} leaks planted in the fitting path. No subset has a false alarm on the protocol (\nv{abl.G1G2G3.fa}\% in every row), so the choice among subsets is a choice among powers. Each gate alone has a blind spot. The first gate alone detects \nv{abl.G1.full} of the \nv{h5.total} leaks in every draw and never detects \nv{abl.G1.never}: the calibration leaks, whose fitted parameters before the calibration map do not depend on the holdout. The second gate alone detects \nv{abl.G2.full} and never detects \nv{abl.G2.never}: the leaks through inputs only, imputation and bin edges for the scorecard and imputation and scaling for the logit, whose fitted values do not depend on labels. The third gate alone detects \nv{abl.G3.full} of the \nv{h5.total} in every draw and misses none altogether, with a mean detection of \nv{abl.G3.mean}\%; what it does not catch in every draw is the choice among configurations made on the holdout (Proposition~\ref{prop:discrete}). The pair of the second and the third gates detects all \nv{abl.G2G3.full} leaks in every draw, and the first gate adds nothing to this pair on the planted leaks. The first gate is nevertheless kept: it localises a failure to a transform instead of the calibration map.

\begin{table}[htbp]
\centering
\small
\caption{Ablation of the gates over the 15 leaks planted in the compact pipeline (100 draws each): the number of leaks that at least one gate of the subset detects in every draw, the number never detected by any gate of the subset, the mean over leaks of the share of draws in which at least one gate of the subset fails, and the largest false alarm rate on the protocol.}
\label{tab:ablation-gates}
\fitwidth{%
\begin{tabular}{lrrrr}
\toprule
Gates kept & Leaks always detected & Leaks never detected & Mean detection \% & False alarm \% \\
\midrule
G1 only & 9 of 15 & 4 & 69.5 & 0 \\
G2 only & 10 of 15 & 4 & 73.3 & 0 \\
G3 only & 13 of 15 & 0 & 95.9 & 0 \\
G1 and G2 & 14 of 15 & 0 & 99.9 & 0 \\
G1 and G3 & 14 of 15 & 0 & 98.0 & 0 \\
G2 and G3 & 15 of 15 & 0 & 100.0 & 0 \\
G1, G2 and G3 & 15 of 15 & 0 & 100.0 & 0 \\
\bottomrule
\end{tabular}%
}
\end{table}


The cost follows from equation~\eqref{eq:cost}: the pair of the second and the third gates needs the base fit, one fit for the flipped labels and $d$ fits for the perturbation draws, and the control one more, so that dropping the first gate saves one fit of the whole path out of $4+d$. The saving is small and the first gate stays in the protocol. The non-vacuity control passes in every draw for every variant of the experiments (\nv{e2.w.none.nv}\% for the scorecard under the protocol), so that none of the passes reported above is a pass of a pipeline that ignores its data; the control is not ablated, since a pipeline that fails it is defective and not leaky.

The ablation also shows what the protocol does not do. It does not choose, among the leaks, those that matter: a gate fails with the same certainty for imputation on all rows, which costs nothing in bias, as for the weight of evidence with holdout labels, which costs \nv{syn.w.woe.bias}. The ranking of the leaks by cost is the one of Table~\ref{tab:main-results}, not the one of Table~\ref{tab:gates}. A validator who finds that the first gate failed has found a breach of the rule and has to read the bias table to know what it costs in the step that failed.
